% Meets and joins of Grothendieck topologies in the sameness lattice.
% Verification: scripts/comparison/topjoin_*.py, tensor_*.py, lemmaP_check.py, monoids_small.py.
\documentclass[11pt]{amsart}
\usepackage{amsmath,amssymb}
\setlength{\emergencystretch}{3em}
\tolerance=600
\usepackage{tikz-cd}
\usepackage{booktabs}
\usepackage[hidelinks]{hyperref}

\newtheorem{thm}{Theorem}[section]
\newtheorem{lem}[thm]{Lemma}
\newtheorem{cor}[thm]{Corollary}
\newtheorem{prop}[thm]{Proposition}
\newtheorem{conj}[thm]{Conjecture}
\theoremstyle{definition}
\newtheorem{defi}[thm]{Definition}
\newtheorem{rem}[thm]{Remark}
\newtheorem{ex}[thm]{Example}

\newcommand{\Same}{\mathrm{Same}}
\newcommand{\Top}{\mathrm{Top}}
\newcommand{\Mor}{\mathrm{Mor}}
\newcommand{\Iso}{\mathrm{Iso}}
\newcommand{\Presh}{\mathrm{Presh}}
\newcommand{\Sh}{\mathrm{Sh}}
\newcommand{\Sat}{\mathrm{Sat}}
\newcommand{\Loc}{\mathrm{Loc}}
\newcommand{\Hom}{\mathrm{Hom}}
\newcommand{\M}{M}
\newcommand{\Ex}{\mathrm{Ex}^\infty}

\begin{document}

\title[Meets and joins of topologies in the sameness lattice]{Meets and joins of
  Grothendieck topologies in the sameness lattice}
\author{Ryota Shibusawa}
\thanks{Corresponding author. E-mail: \texttt{r-shibusawa@daiichi-koudai.ac.jp}}
\address{Daiichi Institute of Technology}
\email{r-shibusawa@daiichi-koudai.ac.jp}
\keywords{Grothendieck topology, subtopos, sheafification, local isomorphism,
  two-out-of-three, idempotent ideal, finite monoid, tensor product of bisets,
  pro-limit, $\infty$-topos, modality, truncation}
\subjclass[2020]{18F10 (primary), 18F20, 18A32, 20M12, 18N60 (secondary)}

\begin{abstract}
For a small category $S$, every Grothendieck topology $J$ determines the class
$W_J$ of maps of presheaves inverted by sheafification, an element of the
complete lattice $\Same(\Presh(S))$ of all two-out-of-three classes. We study
how the lattice operations of $\Top(S)$ and of $\Same(\Presh(S))$ interact
under $J\mapsto W_J$. Meets are always preserved: $W_{\bigwedge J_i}=\bigcap
W_{J_i}$, because the sieves that test local injectivity and surjectivity of a
map do not depend on the topology. Joins are preserved up to saturation,
$W_{J\vee J'}=\Sat(W_J\vee W_{J'})$, and we prove that they are preserved on
the nose for every finite acyclic category (finite posets, parallel-arrow
categories), by a one-object-at-a-time replacement, and for every finite
category whose two topologies have an idempotent intersection of ideals ---
all but four of the $4623$ ordered pairs of idempotent ideals of the monoids
of order at most five. The proof for finite monoids is a
manipulation of tensor products of bisets: the tower of alternating plus
constructions is representable by the tensor products $D_k\otimes\cdots\otimes D_1$
of the ideals, and two factorisation lemmas for idempotent ideals of finite
monoids show that these tensor products collapse to $I\otimes I$, the
representing object of sheafification for the join, after four plus steps. Joins
are \emph{not} preserved in general: on the poset $\mathbb N$ the topologies of
the even and the odd points have the degenerate topology as join, whereas the
join of their samenesses is inverted by the limit of the tower, so it misses
the map from the presheaf of final segments to the terminal presheaf. Thus the
sameness lattice retains the pro-limit invariant that the subtopos lattice
forgets. In every $\infty$-topos, finally, the join of a lex modality with the
$n$-truncation is the $n$-truncated modality, and in simplicial spaces the
join of the realization equivalences with the $(n{+}1)$-coskeletal
equivalences has as its saturation exactly the maps whose realization is an
$n$-equivalence: the identification of Postnikov truncations as joins of
sheaf and homotopy equivalences persists at the $\infty$-level up to
saturation, because a constant simplicial space is $(n{+}1)$-coskeletal iff
its value is $n$-truncated. The finite classifications (monoids of order at most five, their idempotent
ideals, the tensor products and the factorisation lemmas) are verified
exhaustively by computation from the definitions, and the categorical
formulas are cross-checked on extensive finite test sets.
\end{abstract}

\maketitle

\section{Introduction}

A \emph{sameness} on a category $C$ is a class $W$ of morphisms containing the
identities, closed under composition and satisfying two-out-of-three; the
samenesses form a complete lattice $\Same(C)$ under inclusion
\cite{Same,Postnikov}. Two classical sources of samenesses are homotopy theory
(weak equivalences) and sheaf theory: for a Grothendieck topology $J$ on a
small category $S$ the class
\[
W_J=\{f\in\Presh(S):\ a_Jf\text{ is an isomorphism}\}
\]
of $J$-local isomorphisms is a saturated element of $\Same(\Presh(S))$, and
$J\mapsto W_J$ is an order-embedding $\Top(S)\hookrightarrow\Same(\Presh(S))$
\cite[Thm.~3.1]{Postnikov}. This paper asks to what extent this embedding is a
lattice homomorphism.

\medskip\noindent\textbf{Results.}
\begin{itemize}
\item[(A)] \emph{Meets are preserved} (Theorem~\ref{thm:meet}): for every family
of topologies, $W_{\bigwedge_iJ_i}=\bigcap_iW_{J_i}$. The reason is elementary:
a map is a $J$-local isomorphism iff certain sieves attached to the map are
$J$-covering, and these sieves do not depend on $J$.
\item[(B)] \emph{Joins are preserved up to saturation}
(Proposition~\ref{prop:sat}): $W_{J\vee J'}=\Sat(W_J\vee W_{J'})$, and
$W_J\vee W_{J'}=W_{J\vee J'}$ iff the alternating sheafification
$X\to a_JX\to a_{J'}a_JX\to\cdots$ reaches a $(J\vee J')$-sheaf in finitely
many steps for every $X$ (Lemma~\ref{lem:reduction}).
\item[(C)] \emph{Joins are preserved on finite categories with idempotent
intersection} (Theorems~\ref{thm:acyclic} and~\ref{thm:monoid}). For a finite
category, topologies are idempotent two-sided ideals of morphisms
\cite[Thm.~5.1]{Postnikov}, and the join of two topologies is the largest
idempotent ideal inside the intersection of their ideals. If that intersection
is itself idempotent --- the case for $4619$ of the $4623$ ordered pairs of
idempotent ideals of monoids of order $\le5$ --- then $W_J\vee W_{J'}=W_{J\vee J'}$;
and for acyclic finite categories the join is preserved without any
hypothesis, although there the intersection of two ideals is typically
\emph{not} idempotent (already for the two-element chain). For acyclic
categories the proof is a one-object-at-a-time replacement; for monoids it is
a computation with
tensor products of bisets: the tower of alternating plus constructions is represented by
$D_k\otimes_M\cdots\otimes_MD_1$, and two factorisation lemmas for idempotent
ideals of finite monoids (every element factors through an idempotent of the
ideal, with the idempotent absorbed on either side) show that these tensor
products collapse to $I\otimes_MI$ after four plus steps.
\item[(D)] \emph{Joins are not preserved in general} (Theorem~\ref{thm:N}). On
the poset $\mathbb N$, the rigid topologies of the even and of the odd points
have the degenerate topology as their join, so $W_{J\vee J'}$ is everything;
but the limit of the tower inverts both samenesses, hence their join, and does
not invert $X\to1$ for the presheaf of final segments. The sameness lattice
remembers a pro-limit invariant that the subtopos lattice forgets.
\item[(E)] \emph{The $\infty$-categorical manifestation}
(Theorems~\ref{thm:infty} and~\ref{thm:sspace}): in every $\infty$-topos the
join of the sameness of a lex modality $L$ with that of the $n$-truncation is
the sameness of $\tau_{\le n}L$ --- an instance of the join mechanism, since
lex localizations commute with truncation; and in simplicial spaces the
saturation of the join of the realization equivalences with the
$(n{+}1)$-coskeletal equivalences is the class of maps whose realization is an
$n$-equivalence, because a constant simplicial space is $(n{+}1)$-coskeletal
iff its value is $n$-truncated. Whether that join is already saturated, as it
is for simplicial sets \cite{Postnikov}, is left open.
\end{itemize}
The finite computations are described in \S\ref{sec:comp}; the residual case
of (C) for monoids, where the intersection of the two ideals is not
idempotent, is isolated as a single combinatorial statement
(Conjecture~\ref{conj:threefold}), verified on all triples of idempotent ideals
of monoids of order $\le5$; its profunctor form would settle all finite
categories.

\medskip\noindent\textbf{Related work.} Sheafification, local isomorphisms and
the lattice of subtoposes are classical \cite{MLM,Elephant}; the
correspondence between topologies on a finite category and idempotent
two-sided ideals is the finite-sieve form of the Kelly--Lawvere description of
essential localizations \cite{KL89}, and the rigidity of topologies on Cauchy-complete categories with finite
slices, in particular on finite Cauchy-complete categories, is in
\cite{Elephant,Marques25}. Tensor
products of acts over semigroups and the ``firm'' acts of semigroup Morita
theory \cite{LMR18} are the natural home of our bimodule computations; the
factorisation lemmas of \S\ref{sec:monoid} are of the same flavour as the
firmification of acts over factorisable semigroups. Anel, Biedermann, Finster and Joyal \cite{ABFJ} describe left-exact
localizations of $\infty$-topoi by the classes of maps they invert, the
\emph{congruences} (strongly saturated classes with finite-limit and
base-change stability), and study the lattice of these; there joins are taken
at the saturated level. The question here is complementary: whether the join
taken \emph{before} saturation, in the lattice of bare two-out-of-three
classes, already agrees with the saturated one --- Proposition~\ref{prop:sat}
is exactly the comparison of the two joins, and Theorems~\ref{thm:acyclic},
\ref{thm:monoid}, \ref{thm:cat} and~\ref{thm:N} decide it in the cases
considered.

\section{Preliminaries}\label{sec:prelim}

\subsection{Samenesses and saturation}
For a category $C$ let $\Same(C)$ be the complete lattice of samenesses. For
$W\in\Same(C)$ an object $Z$ is \emph{$W$-local} if $\Hom(f,Z)$ is bijective
for all $f\in W$; $\Loc(W)$ is the class of $W$-local objects and
$\Sat(W)=\{f:\Hom(f,Z)\text{ bijective for all }Z\in\Loc(W)\}$ the
saturation. If $L$ is a reflector onto a full replete subcategory
$\mathcal E$, then $W_L=L^*(\Iso)$ is saturated with $\Loc(W_L)=\mathcal E$
\cite[Lemma~2.3]{Postnikov}. We use throughout:

\begin{lem}\label{lem:loc}
$\Loc(U\vee V)=\Loc(U)\cap\Loc(V)$ for $U,V\in\Same(C)$.
\end{lem}
\begin{proof}
The class $\{f:\Hom(f,Z)\text{ bijective}\}$ is a sameness; if it contains
$U$ and $V$ it contains their join.
\end{proof}

\subsection{Topologies, sheafification, the finite dictionary}
For a small category $S$, $\Top(S)$ denotes the complete lattice of
Grothendieck topologies; meets are intersections. $a_J$ is sheafification, the
left-exact reflector onto $\Sh_J(S)$, and $W_J:=a_J^*(\Iso)$ is the class of
$J$-local isomorphisms: $f\colon X\to Y$ is $J$-locally surjective if for every
$c$ and $y\in Y(c)$ the sieve $\{h\colon d\to c: y\cdot h\in\operatorname{im}f_d\}$ is
$J$-covering, and $J$-locally injective if for all $x,x'\in X(c)$ with
$f(x)=f(x')$ the sieve $\{h: x\cdot h=x'\cdot h\}$ is $J$-covering
\cite[III.7]{MLM}.

If every object of $S$ has finitely many sieves (e.g.\ $S$ finite),
topologies correspond, by reverse inclusion, to the \emph{idempotent two-sided
ideals} $D$ of $S$ (classes of morphisms closed under pre- and post-composition
with $D\circ D=D$): $J_D(c)=\{R\supseteq D(c)\}$ and $D_J(c)=\bigcap J(c)$
\cite[Thm.~5.1]{Postnikov}. Under this dictionary $J_D\le J_E$ iff $D\supseteq E$,
$J_D\wedge J_E=J_{D\cup E}$, and $J_D\vee J_E=J_{I}$ where $I$ is the largest
idempotent ideal contained in $D\cap E$ (a union of idempotent ideals is
idempotent, so this largest ideal exists). We write $W_D$ for $W_{J_D}$.

\subsection{The plus construction and local isomorphisms}
For a covering sieve $R$ on $c$ and a presheaf $X$, $X^{+_R}(c)=\Hom(R,X)$;
when $J=J_D$ has least covering sieves the plus construction is $X^+=\Hom(D(-),X)$
and $a_JX=X^{++}$. The unit $X\to X^+$ is a $J$-local isomorphism
\cite[III.5]{MLM}, hence lies in $W_J$.

\section{Meets are preserved}

\begin{thm}\label{thm:meet}
For any small category $S$ and any family $(J_i)_{i\in I}$ of topologies,
$W_{\bigwedge_iJ_i}=\bigcap_iW_{J_i}$. Hence $J\mapsto W_J$ preserves arbitrary
meets, and $\Top(S)$ is a meet-subsemilattice of $\Same(\Presh(S))$.
\end{thm}
\begin{proof}
The meet of topologies is their intersection: a sieve is $(\bigwedge J_i)$-covering
iff it is $J_i$-covering for every $i$. A map $f$ is a $J$-local isomorphism
iff the sieves $\{h:y\cdot h\in\operatorname{im}f\}$ and $\{h:x\cdot h=x'\cdot h\}$
attached to $f$ are $J$-covering; these sieves depend on $f$ only. So $f$ is a
$(\bigwedge J_i)$-local isomorphism iff it is a $J_i$-local isomorphism for all $i$.
\end{proof}

\section{Joins: saturation and the reduction lemma}\label{sec:join}

\begin{prop}\label{prop:sat}
For topologies $J,J'$: $W_{J\vee J'}=\Sat(W_J\vee W_{J'})$. In particular
$W_J\vee W_{J'}\subseteq W_{J\vee J'}$, with equality iff $W_J\vee W_{J'}$ is
saturated.
\end{prop}
\begin{proof}
A presheaf is a $(J\vee J')$-sheaf iff it is a $J$-sheaf and a $J'$-sheaf,
because the sieves for which a given presheaf and all its pullbacks satisfy
the sheaf condition form a topology \cite[C2.1]{Elephant}. By
Lemma~\ref{lem:loc}, $\Loc(W_J\vee W_{J'})=\Sh_J\cap\Sh_{J'}=\Sh_{J\vee J'}
=\Loc(W_{J\vee J'})$, and $W_{J\vee J'}$ is saturated.
\end{proof}

Write $V:=W_J\vee W_{J'}$. Two towers will be used: the \emph{alternating
sheafification tower} $X\to a_JX\to a_{J'}a_JX\to\cdots$ and, when $J,J'$ have
least covering sieves, the \emph{alternating plus tower}
$X\to X^{+_J}\to(X^{+_J})^{+_{J'}}\to\cdots$; in either case write $\Phi_kX$ for
the $k$-th stage. Every map in either tower is a unit of $a_J$ or of $a_{J'}$,
or a plus-construction map $Y\to Y^{+}$, hence lies in $W_J$ or in $W_{J'}$
(\S\ref{sec:prelim}). The lemma below applies verbatim to both towers.

\begin{lem}[reduction]\label{lem:reduction}
\begin{enumerate}
\item For every map $f$: $f\in V\iff a_Jf\in V\iff\Phi_kf\in V$.
\item If $a_{J'}a_Jf$ is an isomorphism then $f=g\circ h$ with $h\in W_J$ and
$g\in W_{J'}$; in particular $f\in V$.
\item If, for one of the two towers, $\Phi_kX$ is a $(J\vee J')$-sheaf for
every $X$ and some $k=k(X)$, then $W_J\vee W_{J'}=W_{J\vee J'}$.
\end{enumerate}
\end{lem}
\begin{proof}
(1) Naturality $\eta_{X'}f=a_J(f)\eta_X$ with $\eta\in W_J\subseteq V$ and
two-out-of-three, twice; the same for each stage of either tower. (2) Let $Z:=X'\times_{a_JX'}a_JX$, with $f=g\circ h$
for the induced $h\colon X\to Z$ and the projection $g\colon Z\to X'$. Then $g$ is
a pullback of $a_Jf$, which is a $J'$-local isomorphism when $a_{J'}a_Jf$ is
iso, and local isomorphisms are pullback-stable; and $Z\to a_JX$ is a pullback
of $\eta_{X'}\in W_J$, so $(Z\to a_JX)\circ h=\eta_X\in W_J$ gives $h\in W_J$.
(3) For $f\colon X\to X'$ in $W_{J\vee J'}$ take $k$ working for $X$ and $X'$;
$\Phi_kf$ is a map between $(J\vee J')$-sheaves inverted by $a_{J\vee J'}$,
hence an isomorphism, so $f\in V$ by (1). See Figure~\ref{fig:reduction}.
\end{proof}

\begin{figure}[ht]
\centering
\begin{tikzcd}[column sep=large,row sep=large]
X\arrow[r,"f"]\arrow[d,"\eta_X"',"W_J"] & X'\arrow[d,"\eta_{X'}","W_J"'] \\
a_JX\arrow[r,"a_Jf"']\arrow[d,"W_{J'}"'] & a_JX'\arrow[d,"W_{J'}"] \\
\Phi_2X\arrow[r,"\Phi_2f"'] & \Phi_2X'
\end{tikzcd}
\qquad
\begin{tikzcd}[column sep=large,row sep=large]
X\arrow[dr,"h","W_J"']\arrow[drr,bend left=18,"f"]\arrow[ddr,bend right=18,"\eta_X"'] & & \\
& Z\arrow[r,"g","W_{J'}"']\arrow[d] & X'\arrow[d,"\eta_{X'}"] \\
& a_JX\arrow[r,"a_Jf"'] & a_JX'
\end{tikzcd}
\caption{Lemma~\ref{lem:reduction}: the alternating tower (left) moves the
question from $f$ to $\Phi_kf$; the pullback trick (right) factors $f$ through
a $W_J$-map and a $W_{J'}$-map when $a_{J'}a_Jf$ is invertible.}
\label{fig:reduction}
\par\smallskip\noindent\textit{Alt text:} Left, a ladder of two commutative
squares from $f$ down to its double sheafification, with vertical units
labelled by the two samenesses. Right, a pullback square of the sheafified map
along the unit of the target, with $f$ factored through the pullback object.
\end{figure}

\section{Finite acyclic categories}\label{sec:acyclic}

Call a finite category \emph{acyclic} if $\Hom(a,b)\neq\varnothing\neq\Hom(b,a)$
implies $a=b$ and the only endomorphisms are identities. Such a category is
Cauchy complete (its only idempotents are identities), every topology on it is
rigid \cite[C2.2.21]{Elephant}, \cite{Marques25}, and by the dictionary the topologies are the subsets
$Y\subseteq\operatorname{Ob}S$: $J_Y$ has ideal $D_Y=\{u\colon a\to b:\ u\text{
factors through some }y\in Y\}$, $\Sh_{J_Y}\simeq\Presh(Y)$ by restriction,
$a_{J_Y}=i_*i^*$ for $i\colon Y\hookrightarrow S$, and
$W_Y=\{f: f_y\text{ bijective for all }y\in Y\}$. The join of $J_Y$ and $J_{Y'}$
is $J_{Y\cap Y'}$: an idempotent ideal of an acyclic finite category is
generated by the identities it contains (an arrow in it factors through
members, and iterating, acyclicity and finiteness force an identity to
appear), so the largest idempotent ideal inside $D_Y\cap D_{Y'}$ is
$D_{Y\cap Y'}$. The intersection $D_Y\cap D_{Y'}$ itself is in general not
idempotent: for the chain $a<b$, $D_{\{a\}}\cap D_{\{b\}}=\{a\to b\}$, and the
two ideals do not commute. So the acyclic case lies in the regime that the
tensor argument of \S\ref{sec:monoid} does not cover, and needs its own proof.

\begin{thm}\label{thm:acyclic}
For a finite acyclic category $S$ and all $Y,Y'\subseteq\operatorname{Ob}S$,
$W_Y\vee W_{Y'}=W_{Y\cap Y'}$. Hence $J\mapsto W_J$ is a lattice embedding
$\Top(S)\cong(2^{\operatorname{Ob}S})^{\mathrm{op}}\hookrightarrow\Same(\Presh(S))$.
\end{thm}
\begin{proof}
Let $f\colon X\to X'$ be bijective on $Y\cap Y'$. Choose a linear extension
$p_1,\dots,p_m$ of the preorder $a\preceq b\iff\Hom(a,b)\neq\varnothing$
(predecessors first) and put $S_i=\{p_1,\dots,p_i\}$, a set closed under
predecessors. Define $X_i(q)=X'(q)$ for $q\in S_i$ and $X(q)$ otherwise, with
restriction along $u\colon q\to q'$ given by $X'(u)$ if both objects are in
$S_i$, by $X(u)$ if both are outside, and by $f_q\circ X(u)$ if $q\in S_i$,
$q'\notin S_i$ (the fourth case cannot occur). Functoriality is naturality of
$f$. The maps $\varphi_i\colon X_{i-1}\to X_i$ (identity away from $p_i$,
$f_{p_i}$ at $p_i$) are natural, $f=\varphi_m\circ\cdots\circ\varphi_1$, and
$\varphi_i$ is bijective away from $p_i$. If $p_i\notin Y$ then
$\varphi_i\in W_Y$; if $p_i\notin Y'$ then $\varphi_i\in W_{Y'}$; and if
$p_i\in Y\cap Y'$ then $\varphi_i$ is an isomorphism. Hence $f\in W_Y\vee W_{Y'}$.
The reverse inclusion is Proposition~\ref{prop:sat}, and meets are
Theorem~\ref{thm:meet} (with $J_Y\wedge J_{Y'}=J_{Y\cup Y'}$).
\end{proof}

\begin{figure}[ht]
\centering
\begin{tikzcd}[column sep=large,row sep=large]
X(q')\arrow[r,"X(u)"]\arrow[d,equal] & X(q)\arrow[d,"f_q"] & & X(p_i)\arrow[r,"f_{p_i}"]\arrow[d,"X(u)"'] & X'(p_i)\arrow[d,"X'(u)"] \\
X_i(q')\arrow[r,"f_q\circ X(u)"'] & X'(q)=X_i(q) & & X(q)\arrow[r,"f_q"'] & X'(q)
\end{tikzcd}
\caption{Theorem~\ref{thm:acyclic}: the mixed presheaf $X_i$ (left: the
restriction from an unprocessed object $q'$ to a processed object $q$ is
$f_q\circ X(u)$) and the naturality square making $\varphi_i$ a map of
presheaves at the newly processed object $p_i$ (right).}
\label{fig:acyclic}
\par\smallskip\noindent\textit{Alt text:} Two commutative squares. Left: the
value of $X$ at an unprocessed object restricts to the value of $X$ at a
processed object and then maps by $f$ to the value of $X'$ there. Right: the
component of $f$ at the newly processed object commutes with the restriction
maps of $X$ and $X'$.
\end{figure}

The acyclicity is what allows the value at a single object to be replaced;
with cycles, an object in $D\setminus D'$ and one in $D'\setminus D$ may be
mutually reachable, and the replacement must be organised differently
(\S\ref{sec:monoid}).

\section{Finite monoids: the tensor tower}\label{sec:monoid}

Let $M$ be a finite monoid, $\Presh(M)$ the category of right $M$-sets, and
$D,E$ idempotent two-sided ideals, with $J_D$, $J_E$ the corresponding
topologies. Right ideals are the sieves; $D$ and $E$ are the least covering
sieves. For a right $M$-set $Z$ and an $M$--$M$-biset $P$ we have the tensor
$Z\otimes_MP$ (quotient of $Z\times P$ by $(zm,p)\sim(z,mp)$, a right $M$-set)
and the adjunction
\begin{equation}\label{eq:adj}
\Hom_M(Z\otimes_MP,X)\cong\Hom_M(Z,\Hom_M(P,X)),
\end{equation}
where $\Hom_M(P,X)$ carries the right action $(\varphi\cdot m)(p)=\varphi(mp)$.
Consequently the alternating plus tower (one plus construction per
factor; sheafification is two consecutive plus constructions) is representable:
\[
\Phi_kX:=\bigl((X^{+_{D}})^{+_{E}}\bigr)^{+_D\cdots}=\Hom_M(T_k,X),
\]
\[
T_k:=D_k\otimes_M\cdots\otimes_MD_1,\qquad D_1=D,\ D_2=E,\ D_3=D,\ \dots,
\]
with the unit $\Phi_{k-1}X\to\Phi_kX$ induced by the multiplication
$\mu_k\colon T_k=D_k\otimes T_{k-1}\to T_{k-1}$. Each unit is in $W_D$ or
$W_E$. For an idempotent ideal $I$, $a_IX=\Hom_M(I\otimes_MI,X)$.

We write $[x_k,\dots,x_1]$ for the class of a tuple in $T_k$; the relations
are $[\dots,am,b,\dots]=[\dots,a,mb,\dots]$ across each tensor sign, subject
to every entry staying in its slot ideal.

\subsection{Two factorisation facts}
\begin{lem}\label{lem:F}
Let $Y$ be an idempotent two-sided ideal of a finite monoid. Every $y\in Y$
admits factorisations
\begin{align*}
\text{(F1)}\qquad& y=y'gy''\ \text{ with }g=g^2\in Y,\ y'g=y',\ y',y''\in Y;\\
\text{(F2)}\qquad& y=y'gy''\ \text{ with }g=g^2\in Y,\ gy''=y'',\ y',y''\in Y.
\end{align*}
\end{lem}
\begin{proof}
Write $y=y_1\cdots y_N$ with $y_i\in Y$ and $N>|M|+2$. Among the prefixes
$p_j=y_1\cdots y_j$ ($1\le j\le|M|+1$) two coincide, $p_i=p_j$ with $i<j$, so
$p_i=p_im$ for $m=y_{i+1}\cdots y_j\in Y$, hence $p_i=p_im^r$ for every $r\ge1$;
since $M$ is finite some power $g=m^r$ is idempotent, $g\in Y$ because $Y$ is
closed under products, and $p_i=p_ig$. Then $y=p_i\,g\,(y_{j+1}\cdots y_N)$ is (F1).
(F2) is the same argument on suffixes.
\end{proof}

\subsection{Merging lemmas}
\begin{lem}[P]\label{lem:P}
Let $X,Y,Z$ be two-sided ideals with $Y$ idempotent. In $X\otimes Y\otimes Z$,
for $c\in Y\cdot Z$ the class $[x,y,c]$ depends only on $(xy,c)$.
\end{lem}
\begin{proof}
Fix a rule writing $c=y_1z_1$ ($y_1\in Y$, $z_1\in Z$) and $y_1=y'gy''$ by (F1).
Moving $y'g$ to the left (the remaining $y''z_1$ lies in $Z$) and then $yy'$ to
the left (the remaining $g$ lies in $Y$):
\[
[x,y,c]=[x,\,y\,y'g,\,y''z_1]=[xyy',\,g,\,y''z_1],
\]
which depends on $xy$ and on the data $(y',g,y'',z_1)$ determined by $c$.
\end{proof}
\begin{lem}[P$'$]\label{lem:Pp}
With $Y$ idempotent, for $c\in X\cdot Y$ the class $[c,y,z]\in X\otimes Y\otimes Z$
depends only on $(c,yz)$. (Mirror image, using (F2).)
\end{lem}

\begin{lem}[absorption]\label{lem:absorb}
If $D_j$ is idempotent, $[\dots,x_j,x_{j-1},\dots]=[\dots,x_j',x_j''x_{j-1},\dots]$
for any factorisation $x_j=x_j'x_j''$ in $D_j$; so an interior slot may be
assumed to lie in the product of its ideal with a neighbouring ideal.
\end{lem}

\begin{lem}[merging]\label{lem:merge}
Let $X\otimes Y\otimes Z\otimes\cdots$ be a tensor of ideals with $Y$ idempotent.
If $[x,y,c,\dots]$ depends only on $(xy,c,\dots)$ for \emph{all} tuples, then
the multiplication $X\otimes Y\otimes Z\otimes\cdots\to XY\otimes Z\otimes\cdots$
is a bijection (and likewise for merging a later pair under the hypothesis of
Lemma~\textup{P$'$}).
\end{lem}
\begin{proof}
The map is surjective. Define a section on tuples of the target by
$L[w,c,\dots]:=[x,y,c,\dots]$ for any factorisation $w=xy$; it is well defined
by hypothesis, and it respects the relations of the target: for $[wm,c]\sim[w,mc]$
use the factorisation $wm=x(ym)$; for $[w,cm]\sim[w,c,m\cdot]$ nothing changes
in the first two slots. So $L$ descends to the target, and $L\circ\mu=\mathrm{id}$.
\end{proof}

\subsection{The theorem}
\begin{thm}\label{thm:monoid}
Let $M$ be a finite monoid and $D,E$ idempotent two-sided ideals such that
$I:=D\cap E$ is idempotent (so $J_D\vee J_E=J_I$). Then the canonical
multiplication maps give isomorphisms of bisets $T_k\cong I\otimes_MI$ for all
$k\ge4$. Consequently $\Phi_4X\cong a_IX$ for every $M$-set $X$: the alternating
plus tower reaches the $J_I$-sheafification after four plus steps, and
\[
W_D\vee W_E=W_{D\cap E}\quad\text{in }\Same(\Presh(M)).
\]
\end{thm}
\begin{proof}
$DE\subseteq I$ and $ED\subseteq I$, $I\subseteq D$, $I\subseteq E$, and every
product of two or more alternating factors $D,E$ lies between $I=I^2$ and
$D\cap E=I$, hence equals $I$. Write $T_k=[x_1,\dots,x_k]$ with slot ideals
$D_k,D_{k-1},\dots,D_1$ from left to right.

\emph{Step 1: merge slots $1,2$.} Let $t,t'$ be tuples differing only in slots
$1,2$, with equal products $x_1x_2=x_1'x_2'$. Absorb slot $4$ into slot $3$
in both, with the same factorisation $x_4=pq$ (Lemma~\ref{lem:absorb}; here
$k\ge4$ is used): slot $3$ becomes $x_3p\in D_{k-2}D_{k-3}=I\subseteq
D_{k-1}D_{k-2}$. Lemma~P for the triple of slots $1,2,3$ shows that the two
classes agree. By Lemma~\ref{lem:merge} the multiplication
$T_k\to(D_kD_{k-1})\otimes D_{k-2}\otimes\cdots\otimes D_1=I\otimes D_{k-2}\otimes
\cdots\otimes D_1$ is a bijection.

\emph{Step 2: merge slots $2,3$ repeatedly.} In $I\otimes Y\otimes Z\otimes\cdots$
with $Y$ idempotent, every element $c$ of the first slot lies in $I=I\cdot I
\subseteq I\cdot Y$, so Lemma~P$'$ applies to every tuple and
Lemma~\ref{lem:merge} gives a bijection onto $I\otimes YZ\otimes\cdots$. Starting
from $I\otimes D_{k-2}\otimes D_{k-3}\otimes\cdots\otimes D_1$, the first such
merge produces $I\otimes I\otimes D_{k-4}\otimes\cdots$ ($D_{k-2}D_{k-3}=I$), and
since $I$ is idempotent the merge can be repeated ($I\cdot D_j=I$) until two
slots remain: $T_k\cong I\otimes I$. All maps are multiplications, so the
composite $T_k\to I\otimes I$ is the canonical one and is an isomorphism of
$M$--$M$-bisets (Figure~\ref{fig:tower}).

Therefore $\Phi_kX=\Hom_M(T_k,X)\cong\Hom_M(I\otimes I,X)=a_IX$ as $M$-sets, so
$\Phi_kX$ is an $I$-sheaf; and $X\to\Phi_kX$ is a composite of units in
$W_D\cup W_E\subseteq W_I$, so $\Phi_kX\cong a_IX$ under $X$ and the unit
$\eta_X\colon X\to a_IX$ lies in $W_D\vee W_E$. For $f\in W_I$ the naturality
square $\eta_{X'}f=a_I(f)\eta_X$ with $a_If$ invertible gives $f\in W_D\vee W_E$.
The reverse inclusion is Proposition~\ref{prop:sat}.
\end{proof}

\begin{figure}[ht]
\centering
\begin{tikzcd}[column sep=huge,row sep=large]
E\otimes D\otimes E\otimes D\arrow[r,"{\text{P},\ (1,2)}"] & I\otimes E\otimes D\arrow[r,"{\text{P}',\ (2,3)}"] & I\otimes I
\end{tikzcd}
\caption{The collapse of $T_4$ in Theorem~\ref{thm:monoid}: the first merge
uses Lemma~P after absorbing the last slot into the third, the second uses
Lemma~P$'$ with the first slot already in $I$. Both arrows are multiplication
maps and both are bijections.}
\label{fig:tower}
\par\smallskip\noindent\textit{Alt text:} A chain of two multiplication
maps: the four-fold tensor of $E$, $D$, $E$, $D$ merges its first pair into $I$,
then the three-fold tensor merges its last pair, ending at $I$ tensor $I$.
\end{figure}

\subsection{Finite categories}\label{sub:cat}
Let $S$ be a finite category. A \emph{two-sided ideal} $D$ of $S$ is a class of
morphisms closed under pre- and post-composition, i.e.\ a sub-profunctor
$D\subseteq\Hom_S$; it is \emph{idempotent} if $D\circ D=D$ (every member
factors as a composite of two members). For ideals (or, more generally,
sub-profunctors of $\Hom_S$) $P,Q$ the composite profunctor is
\[
(Q\otimes_SP)(a,c)=\Bigl(\coprod_{b}Q(b,c)\times P(a,b)\Bigr)\Big/\bigl((q\circ m,\,p)\sim(q,\,m\circ p)\bigr),
\]
an $S$--$S$-profunctor, and for a presheaf $X$ and a profunctor $P$ the presheaf
$\Hom(P,X)(c):=\Hom_{\Presh(S)}(P(-,c),X)$ satisfies
$\Hom(Q\otimes_SP,X)\cong\Hom(Q,\Hom(P,X))$, the profunctor form of
\eqref{eq:adj}. If $J=J_D$ then $X^{+}=\Hom(D,X)$, $a_JX=X^{++}$, and the
alternating plus tower is $\Phi_kX=\Hom(T_k,X)$ with $T_k=D_k\otimes_S\cdots\otimes_SD_1$.
Elements of $T_k(a,c)$ are classes $[x_k,\dots,x_1]$ of composable chains
$a\xrightarrow{x_1}\cdots\xrightarrow{x_k}c$ with $x_i\in D_i$, modulo the
relations $[\dots,x_{i+1}\circ m,x_i,\dots]=[\dots,x_{i+1},m\circ x_i,\dots]$
with all entries in their slot ideals.

\begin{lem}\label{lem:Fcat}
Let $Y$ be an idempotent ideal of the finite category $S$. Every $y\in Y(a,c)$
admits factorisations of both types: $y=y'\circ g\circ y''$ with $g\in Y$ an
idempotent endomorphism, $y',y''\in Y$, and $y'\circ g=y'$ \textup{(F1)}; and
likewise one with $g\circ y''=y''$ \textup{(F2)}.
\end{lem}
\begin{proof}
Write $y=y_N\circ\cdots\circ y_1$ with $y_i\in Y$ and $N$ larger than the number
of morphisms of $S$ plus two. The partial composites $p_j=y_j\circ\cdots\circ y_1$
all have domain $a$, so two coincide, $p_i=p_j$ with $i<j$; then $p_i=m\circ p_i$
for the endomorphism $m=y_j\circ\cdots\circ y_{i+1}\in Y$ of $\operatorname{cod}(p_i)$,
some power $g=m^r$ is idempotent, $g\in Y$, and $g\circ p_i=p_i$. Then
$y=(y_N\circ\cdots\circ y_{j+1})\circ g\circ p_i$ is (F2); (F1) is the same
argument with the partial composites $y_N\circ\cdots\circ y_j$ of common
codomain $c$.
\end{proof}

\begin{thm}\label{thm:cat}
Let $S$ be a finite category and $J_D,J_E$ topologies whose ideals have an
idempotent intersection $I=D\cap E$ (so $J_D\vee J_E=J_I$). Then
$T_k\cong I\otimes_SI$ canonically for $k\ge4$, and $W_D\vee W_E=W_{D\cap E}$ in
$\Same(\Presh(S))$.
\end{thm}
\begin{proof}
Lemmas~P, P$'$, \ref{lem:absorb} and~\ref{lem:merge} are statements about
classes of composable chains and the relations across tensor signs; their
proofs use only the relations, the closure of ideals under composition on
both sides, and Lemma~\ref{lem:F}, which Lemma~\ref{lem:Fcat} replaces. (In
Lemma~P, for $c\in(Y\circ Z)(a,b)$ write $c=y_1\circ z_1$ and $y_1=y'\circ g\circ y''$
by (F1); then $[x,y,c]=[x,\,y\circ y'\circ g,\,y''\circ z_1]=[x\circ y\circ y',\,g,\,y''\circ z_1]$,
each step legitimate because the remaining entry lies in its slot ideal.)
The proof of Theorem~\ref{thm:monoid} is then repeated word for word: after
absorbing slot $4$ into slot $3$, Lemma~P merges slots $1,2$ into
$D_kD_{k-1}=I$; Lemma~P$'$ then merges slots $2,3$ repeatedly, since the first
slot lies in $I=I\circ I\subseteq I\circ Y$; every alternating product of
$\ge2$ factors equals $I$. Hence $\Phi_kX=\Hom(T_k,X)\cong\Hom(I\otimes I,X)=a_IX$
for $k\ge4$, and the last paragraph of the proof of Theorem~\ref{thm:monoid}
gives $W_D\vee W_E=W_I$.
\end{proof}

Theorem~\ref{thm:acyclic} is \emph{not} a special case: for acyclic $S$ the
intersection of two ideals is usually not idempotent, and the replacement
argument of \S\ref{sec:acyclic} handles that regime directly. The two
theorems overlap (acyclic categories with idempotent intersections) but
neither contains the other.

\begin{ex}[semilattice monoids]
For a finite meet-semilattice $L$ as a monoid, ideals are down-sets and all
are idempotent, so Theorem~\ref{thm:monoid} applies to every pair; for
principal down-sets $D={\downarrow}e$, $E={\downarrow}e'$ the unit chain is
explicit, $X\to X\cdot e\to X\cdot e\cdot e'=X\cdot(e\wedge e')$, two steps.
\end{ex}

\subsection{The residual case}
The hypothesis of Theorem~\ref{thm:monoid} can fail: among the $4623$ ordered pairs of
idempotent ideals of the $273$ monoids of order $\le5$ there are four pairs
(in two monoids of order $5$) with $D\cap E$ not idempotent; there $DE\neq ED$
and the join corresponds to the smaller ideal $I'=(DE)^m$. In these cases
the alternating tower still stabilises, to $I'\otimes I'$, in at most four
steps, and the collapse is explicit (an element $2\in D\cap E$ with $2=3\cdot2
=2\cdot4$, $3\in D$ and $4\in E$ idempotent, $2\cdot3=0$). The missing
ingredient is the following statement. The unrestricted version (allowing
$Z=M$) fails in exactly six of the $10\,235$ triples of idempotent ideals of
monoids of order $\le5$, all with $Z=M$, where it reduces to the false
two-fold statement $X\otimes Y\to XY$; under the hypothesis $Z\neq M$ no
failure occurs.

\begin{conj}[three-fold multiplicativity]\label{conj:threefold}
For idempotent two-sided ideals $X,Y,Z\neq M$ of a finite monoid, the class
$[x,y,c]\in X\otimes Y\otimes Z$ depends only on $(xy,c)$ for \emph{all}
$c\in Z$; equivalently $X\otimes Y\otimes Z\to XY\otimes Z$ is a bijection.
\end{conj}
By Lemma~\ref{lem:merge} the conjecture implies $T_k\cong I'\otimes I'$ for
$k$ large and hence $W_D\vee W_E=W_{J_D\vee J_E}$ for all finite monoids. Note that the non-commuting regime $DE\neq ED$ is not exotic: it is the
generic situation for acyclic categories (\S\ref{sec:acyclic}), where join
preservation nevertheless holds by the replacement argument; a proof for
finite categories in general would have to combine the two mechanisms. For a
finite category the corresponding statement is the profunctor form of the
conjecture --- for idempotent ideals $X,Y,Z$ of $S$ with $Z$ not containing any
identity, the class $[x,y,c]\in X\otimes_SY\otimes_SZ$ depends only on
$(x\circ y,c)$ --- which does not follow from the monoid case and is left as a
separate question; with it, the argument of \S\ref{sub:cat} would give join
preservation for all finite categories.

\section{Joins are not preserved in general}\label{sec:N}

\begin{thm}\label{thm:N}
Let $S=(\mathbb N,\le)$ and let $J_{\mathrm{ev}}$, $J_{\mathrm{od}}$ be the
topologies of the ideals $D_Y=\{n\to m:\ \exists y\in Y,\ n\le y\le m\}$ for
$Y$ the even and the odd numbers. Then $J_{\mathrm{ev}}\vee J_{\mathrm{od}}$ is
the degenerate topology, so $W_{J_{\mathrm{ev}}\vee J_{\mathrm{od}}}=\Mor$, but
$W_{J_{\mathrm{ev}}}\vee W_{J_{\mathrm{od}}}\subsetneq\Mor$: it is contained in
the class of maps inverted by $\lim_{\mathbb N^{\mathrm{op}}}\colon\Presh(\mathbb N)\to\mathbf{Set}$.
\end{thm}
\begin{proof}
$D_Y$ is an idempotent two-sided ideal, so $J_Y$ is a topology
\cite[Thm.~5.1]{Postnikov} (that direction needs no finiteness); it is rigid
with irreducible objects $Y$, $\Sh_{J_Y}\simeq\Presh(Y)$ by restriction,
$(a_YX)(m)=X(\max\{y\in Y:y\le m\})$ when the maximum exists and
$(a_YX)(m)=1$ otherwise (for $Y$ the odd numbers and $m=0$ the indexing
category is empty), and $W_Y=\{f:f_y\text{ bijective for }y\in Y\}$.
The join $J$ is degenerate: we show $\varnothing\in J(m)$ by induction on
$m$. For $m=0$ the least $J_{\mathrm{od}}$-covering sieve $D_{\mathrm{od}}(0)$ is
empty. For $m>0$ let $R=D_{\mathrm{od}}(m)$ if $m$ is even and $R=D_{\mathrm{ev}}(m)$
if $m$ is odd; $R$ is $J$-covering, and every arrow $n\to m$ in $R$ has
$n\le y<m$ for some $y$ of the other parity, so $n<m$ and $\varnothing\in J(n)$
by induction; transitivity gives $\varnothing\in J(m)$. Now $\Lambda:=\lim_{\mathbb N^{\mathrm{op}}}$ inverts
$W_{\mathrm{ev}}$ and $W_{\mathrm{od}}$, since evens and odds are cofinal in the
tower, hence inverts their join. For $X(n)=\{x\in\mathbb N:x\ge n\}$ with
inclusions, $\Lambda X=\varnothing$, so $\Lambda(X\to1)$ is not bijective and
$X\to1$ is not in the join.
\end{proof}

The mechanism is the failure of objectwise stabilisation in
Lemma~\ref{lem:reduction}: here $(\Phi_kX)(m)\approx X(m-k)$ shifts by one at
each step and never reaches the terminal sheaf. $\Same(\Presh(\mathbb N))$
retains the pro-limit invariant $\Lambda$ that the lattice of subtoposes
forgets; the saturation of the join is everything, the join itself is a
proper sameness.

\begin{figure}[ht]
\centering
\begin{tikzcd}[column sep=normal,row sep=small]
\cdots\arrow[r] & X(3)\arrow[r]\arrow[d,equal] & X(2)\arrow[r]\arrow[d] & X(1)\arrow[r]\arrow[d,equal] & X(0)\arrow[d] \\
\cdots\arrow[r] & X(3)\arrow[r] & X(2)\arrow[r] & X(1)\arrow[r] & \ast
\end{tikzcd}
\caption{Theorem~\ref{thm:N}: a map bijective on the odd points (vertical
equalities) may change the even values arbitrarily; iterating with the even
points shifts the tower by one step at a time, and the limit $\Lambda$ never
changes.}
\label{fig:N}
\par\smallskip\noindent\textit{Alt text:} Two towers of sets indexed by the
natural numbers, connected by vertical maps that are equalities at the odd
indices and arbitrary at the even indices.
\end{figure}

\section{The $\infty$-categorical manifestation}\label{sec:infty}

Samenesses of an $\infty$-category are classes of morphisms containing the
equivalences, closed under composition and two-out-of-three; the join theorem
of \cite[Thm.~4.1]{Postnikov} and its corollary for two reflectors hold
verbatim, as does Lemma~\ref{lem:loc}. The results of this section are
instances of that mechanism in the well-developed theory of localizations
of $\infty$-topoi \cite{HTT}; we record them because they say precisely what
survives of the sheaf/homotopy comparison at the $\infty$-level.

Let $\mathcal E$ be an $\infty$-topos, $L$ the reflector of a subtopos (an
accessible left-exact localization; a lex modality), and $\tau_{\le n}$ the
$n$-truncation reflector.

\begin{thm}\label{thm:infty}
In $\Same(\mathcal E)$: $W_L\vee W_{\tau_{\le n}}=W_{\tau_{\le n}L}=W_{L\tau_{\le n}}$,
the class of maps that become equivalences after passing to the
$n$-truncated $L$-sheaf; its local objects are the $n$-truncated $L$-local
objects, and $\tau_{\le n}L$ is again a reflective subuniverse. If the subtopos
of $L$ is hypercomplete, $\bigcap_n(W_L\vee W_{\tau_{\le n}})=W_L$.
\end{thm}
\begin{proof}
By \cite[Cor.~4.2]{Postnikov} it suffices that $\tau_{\le n}L$ inverts
$W_{\tau_{\le n}}$. A left-exact left adjoint between presentable
$\infty$-categories commutes with truncation, $L\tau_{\le n}\simeq\tau_{\le n}L$
\cite[5.5.6.28]{HTT}, so $\tau_{\le n}Lg\simeq L\tau_{\le n}g$ is an equivalence
whenever $\tau_{\le n}g$ is. The intersection consists of the maps $f$ with
$Lf$ $\infty$-connective, which are the $L$-equivalences when the subtopos is
hypercomplete.
\end{proof}

Now let $\mathcal P=\mathrm{Fun}(\Delta^{\mathrm{op}},\mathcal S)$ be the
$\infty$-category of simplicial spaces, $|-|\colon\mathcal P\to\mathcal S$ the
colimit (geometric realization), left adjoint to the fully faithful constant
functor $c$ (fully faithful because $\Delta^{\mathrm{op}}$ has the initial object
$[0]$, so its nerve is contractible and colimits of constant diagrams are
trivial), and
$\mathrm{cosk}_{n+1}$ the $(n{+}1)$-coskeleton, the reflector onto the
$(n{+}1)$-coskeletal objects. Put $W_{|-|}:=\{f:|f|\text{ an equivalence}\}$,
$W_{\mathrm{cosk}_{n+1}}:=\{f:\mathrm{cosk}_{n+1}f\text{ an equivalence}\}$ and
$W_{\tau_{\le n}|-|}:=\{f:\tau_{\le n}|f|\text{ an equivalence}\}$.

\begin{lem}\label{lem:constcosk}
For $K\in\mathcal S$ and $n\ge0$, the constant simplicial space $cK$ is
$(n{+}1)$-coskeletal iff $K$ is $n$-truncated.
\end{lem}
\begin{proof}
$(\mathrm{cosk}_{n+1}cK)_m=\mathrm{Map}_{\mathcal P}(\mathrm{sk}_{n+1}\Delta^m,cK)
=\mathrm{Map}_{\mathcal S}(|\mathrm{sk}_{n+1}\Delta^m|,K)$, and the unit
$cK\to\mathrm{cosk}_{n+1}cK$ is, at level $m$, the inclusion of the constant
maps. The space $|\mathrm{sk}_{n+1}\Delta^m|$ is contractible for $m\le n+1$ and
$n$-connected for $m\ge n+2$ (the $(n{+}1)$-skeleton of a contractible CW
complex), with $|\mathrm{sk}_{n+1}\Delta^{n+2}|=|\partial\Delta^{n+2}|\simeq S^{n+1}$.
If $K$ is $n$-truncated then $\mathrm{Map}(A,K)\simeq\mathrm{Map}(\tau_{\le n}A,K)
\simeq K$ for every $n$-connected $A$, so $cK$ is $(n{+}1)$-coskeletal.
Conversely, at level $n+2$ the unit is $K\to\mathrm{Map}(S^{n+1},K)$, whose
fibre over the component of $x$ is $\Omega^{n+1}(K,x)$; it is an equivalence
iff $\pi_i(K,x)=0$ for all $i\ge n+1$ and all $x$, i.e.\ iff $K$ is
$n$-truncated.
\end{proof}

\begin{thm}\label{thm:sspace}
In $\Same(\mathcal P)$, for every $n\ge0$,
\[
\Sat\bigl(W_{|-|}\vee W_{\mathrm{cosk}_{n+1}}\bigr)=W_{\tau_{\le n}|-|},
\]
and the local objects of the join are the constant $n$-truncated simplicial
spaces.
\end{thm}
\begin{proof}
$W_{|-|}$ and $W_{\mathrm{cosk}_{n+1}}$ are the classes inverted by the reflectors
$c|-|$ and $\mathrm{cosk}_{n+1}$, so they are saturated with local objects the
constant and the $(n{+}1)$-coskeletal objects respectively. By
Lemma~\ref{lem:loc} and Lemma~\ref{lem:constcosk} the local objects of the join
are the constant $n$-truncated objects $c(\mathcal S_{\le n})$. This is a
reflective subcategory with reflector $c\circ\tau_{\le n}\circ|-|$ (reflect
into constants, then truncate the value), whose inverted class is
$W_{\tau_{\le n}|-|}$; it is saturated with the same local objects, hence equals
the saturation of the join.
\end{proof}

Thus the theorem ``Postnikov truncation $=$ homotopy $\vee$ coskeleton'' of
\cite{Postnikov} persists at the $\infty$-level up to saturation. For
simplicial sets the join is already saturated, the proof there using the Kan
replacement inside the homotopy sameness; whether
$W_{|-|}\vee W_{\mathrm{cosk}_{n+1}}$ is saturated in simplicial spaces we leave
open. (For $n=-1$ the same statements hold with ``$(-1)$-truncated'' meaning
empty or contractible.)

\section{Computational record}\label{sec:comp}
The classifications below (monoids, ideals, tensor products, factorisation
lemmas) are exhaustive within the stated ranges; the tests on presheaves and
$M$-sets are cross-checks on finite test sets, not exhaustive
(\texttt{scripts/comparison/}, Python~3, standard library):
\begin{itemize}
\item \texttt{topjoin\_poset.py}: for all posets on $\le4$ points and $25$ random
finite presheaves each, the alternating sheafification reaches a
$(Y\cap Y')$-sheaf in at most $n+1$ steps ($>110\,000$ cases).
\item \texttt{monoids\_small.py}: backtracking enumeration of monoids of order
$\le5$ ($1,2,7,35,228$, matching the known counts).
\item \texttt{topjoin\_monoid.py}: on the $23$ monoids of order $\le5$ with
incomparable idempotent ideals, the alternating chain stabilises in $\le3$
steps on all cyclic $M$-sets of size $\le4$, and Theorem~\ref{thm:meet} was
checked on $12\,442$ equivariant maps.
\item \texttt{tensor\_tower.py}, \texttt{tensor\_multi.py}: the towers $T_k$ for
all $66$ incomparable pairs stabilise by $k=4$ to the size of $I\otimes I$
(resp.\ $I'\otimes I'$); $E\otimes D\cong E\cap D$ and $I\otimes I\cong I$ fail
in general.
\item \texttt{lemmaP\_check.py}: Lemmas~P and~P$'$ on all $4245$ triples of
proper idempotent ideals; \texttt{tensor\_assoc.py}: the unrestricted form of
Conjecture~\ref{conj:threefold} on all $10\,235$ triples --- six failures, all
with $Z=M$, none under the hypothesis $Z\neq M$.
\end{itemize}

\section{Conclusion}
Under $J\mapsto W_J$ the lattice of topologies sits inside the sameness
lattice as a meet-subsemilattice, and as a sublattice whenever the category is
finite and the two ideals intersect idempotently. The counterexample
exhibited here is genuinely infinite and is detected by a pro-limit; whether
join preservation can fail for a finite category remains open; for finite
monoids it is reduced to a single statement about tensor products of
idempotent ideals (Conjecture~\ref{conj:threefold}), and for finite categories
to its profunctor form.

\section*{Funding}
No funding was received for this work.

\section*{Declarations}
\noindent\textbf{Competing interests.}\quad The author declares that he has no
competing interests.

\noindent\textbf{Availability of data and materials.}\quad All verification
scripts named in \S\ref{sec:comp} are provided with the manuscript.

\begin{thebibliography}{9}
\bibitem{Same} R.~Shibusawa, \emph{The sameness lattice of a category:
weak-equivalence structures, Galois functoriality, and comparison intervals},
preprint (2026), archived at DOI 10.5281/zenodo.22823772.
\bibitem{Postnikov} R.~Shibusawa, \emph{Postnikov truncations as joins of sheaf
and homotopy equivalences}, preprint (2026), archived at DOI 10.5281/zenodo.22956181.
\bibitem{MLM} S.~Mac Lane, I.~Moerdijk, \emph{Sheaves in Geometry and Logic},
Universitext, Springer, 1992.
\bibitem{Elephant} P.~T.~Johnstone, \emph{Sketches of an Elephant: A Topos
Theory Compendium}, Oxford Logic Guides \textbf{43--44}, Oxford University
Press, 2002.
\bibitem{KL89} G.~M.~Kelly, F.~W.~Lawvere, \emph{On the complete lattice of
essential localizations}, Bull. Soc. Math. Belg. S\'er. A \textbf{41} (1989),
289--319.
\bibitem{Marques25} J.~Marqu\`es, \emph{A criterion for categories on which
every Grothendieck topology is rigid}, Appl. Categ. Structures \textbf{33}
(2025), art.\ 37.
\bibitem{LMR18} V.~Laan, L.~M\'arki, \"U.~Reimaa, \emph{Morita equivalence of
semigroups revisited: firm semigroups}, J. Algebra \textbf{505} (2018), 247--270.
\bibitem{ABFJ} M.~Anel, G.~Biedermann, E.~Finster, A.~Joyal, \emph{Left-exact
localizations of $\infty$-topoi I: Higher sheaves}, Adv. Math. \textbf{400}
(2022), 108268.
\bibitem{HTT} J.~Lurie, \emph{Higher Topos Theory}, Annals of Mathematics
Studies \textbf{170}, Princeton University Press, 2009.
\end{thebibliography}
\end{document}
